\section{Classical gauge fluctuation operators on the positive radial cusp}
\label{cusp-gauge-section}

Retain the positive radial flat tori
\(F_T=\mathbb R^4/P_T\mathbb Z^4\), their Euclidean quotient metrics, and
the real quantities \(D_T,L_T,q_T,m_T,B_T,H,B\) from
\cref{cusp-spec-exact-dual-lattice,cusp-spec-shortest-modes}.  Thus
\[
 D_T=L_Tq_T+6m_T^2>0,\qquad
 L_T=\operatorname{Im}\tau_T,\qquad
 q_T=L_T-B_T=-\operatorname{Im}\beta_T,
\]
and \(D_T\) is the positive covolume.  Let
\[
 \mathbf e_\gamma=(1,0,0,0)^{\mathsf t},\qquad
 \mathbf e_u=(0,1,0,0)^{\mathsf t};
\]
then the already-proved dual covectors are
\[
 \ell_{\gamma,T}=P_T^{-\mathsf T}\mathbf e_\gamma,\qquad
 \ell_{u,T}=P_T^{-\mathsf T}\mathbf e_u,
\]
with exact squared lengths
\begin{equation}
 a_{\gamma,T}:=\lVert\ell_{\gamma,T}\rVert^2
 =\frac{L_T^2+m_T^2}{D_T^2},\qquad
 a_{u,T}:=\lVert\ell_{u,T}\rVert^2
 =\frac{q_T^2+36m_T^2}{D_T^2}.
 \label{cusp-gauge-short-mode-lengths}
\end{equation}
For all sufficiently large \(T\), the shortest nonzero dual vectors are
precisely the minimizing members of
\(\{\pm\ell_{\gamma,T},\pm\ell_{u,T}\}\).  Put
\begin{equation}
 G_T=P_T^{\mathsf T}P_T.
 \label{cusp-gauge-coordinate-gram}
\end{equation}
The matrix \(G_T\) is the tangent Gram matrix in the lattice coordinates
\(y\in\mathbb R^4/\mathbb Z^4\), where \(x=P_Ty\).

\subsection{The continuum Yang--Mills Hessian}
\label{cusp-gauge-continuum-subsection}

Let \(G\) be a compact Lie group with real Lie algebra \(\mathfrak g\).
Fix a positive-definite
\(\operatorname{Ad}(G)\)-invariant inner product on \(\mathfrak g\), put
\[
 d_G=\dim_{\mathbb R}\mathfrak g,\qquad
 \mathfrak g_{\mathbb C}=\mathfrak g\otimes_{\mathbb R}\mathbb C,
\]
and fix a coupling \(g_{\mathrm{YM}}>0\).  On the trivial principal bundle
\[
 \mathcal P_T=F_T\times G\longrightarrow F_T
\]
let \(A_0\) be the product connection.  The Yang--Mills functional is
normalized by
\begin{equation}
 \operatorname{YM}_T(A)
 =\frac1{2g_{\mathrm{YM}}^2}
 \int_{F_T}\lvert F_A\rvert^2\,\mathrm d\operatorname{vol}_{F_T}.
 \label{cusp-gauge-ym-normalization}
\end{equation}

\begin{proposition}[Hessian at the product connection]
\label{cusp-gauge-continuum-hessian}
For smooth
\(a,b\in\Omega^1(F_T;\mathfrak g)\),
\begin{equation}
 \operatorname{Hess}_{A_0}\operatorname{YM}_T(a,b)
 =\frac1{g_{\mathrm{YM}}^2}
 \langle\mathrm da,\mathrm db\rangle_{L^2}.
 \label{cusp-gauge-continuum-hessian-form}
\end{equation}
Relative to the unscaled \(L^2\) inner product on one-forms, the associated
closed Hessian form is
\begin{equation}
 \mathfrak l_T[a]
 =g_{\mathrm{YM}}^{-2}\lVert \mathrm d_1a\rVert_{L^2}^2,
 \qquad
 \operatorname{Dom}(\mathfrak l_T)=\operatorname{Dom}(\mathrm d_1),
 \label{cusp-gauge-continuum-closed-form}
\end{equation}
on \(L^2\Omega^1(F_T;\mathfrak g)\), where \(\mathrm d_1\) is the
maximal closed exterior derivative
\[
 \operatorname{Dom}(\mathrm d_1)
 =\{a\in L^2\Omega^1(F_T;\mathfrak g):
       \mathrm da\in L^2\Omega^2(F_T;\mathfrak g)
       \text{ distributionally}\}.
\]
The associated nonnegative self-adjoint operator is
\begin{equation}
 \mathcal L_T=g_{\mathrm{YM}}^{-2}\mathrm d_1^*\mathrm d_1,
 \qquad
 \operatorname{Dom}(\mathcal L_T)
 =\{a\in\operatorname{Dom}(\mathrm d_1):
       \mathrm d_1a\in\operatorname{Dom}(\mathrm d_1^*)\}.
 \label{cusp-gauge-continuum-hessian-operator}
\end{equation}
Its differential expression on smooth one-forms is
\(g_{\mathrm{YM}}^{-2}\mathrm d^*\mathrm d\).
\end{proposition}

\begin{proof}
With the convention
\(F_{A_0+a}=\mathrm da+\tfrac12[a\wedge a]\), one has
\[
 F_{A_0+sa+tb}
 =s\,\mathrm da+t\,\mathrm db
 +O\bigl((|s|+|t|)^2\bigr).
\]
The coefficient of \(st\) in the squared norm is
\(2\langle\mathrm da,\mathrm db\rangle_{L^2}\).  Taking the mixed
derivative in \cref{cusp-gauge-ym-normalization} proves
\cref{cusp-gauge-continuum-hessian-form}.  The maximal distributional
exterior derivative is closed and densely defined.  Hence its graph-norm
square is a closed nonnegative form, and the first representation theorem
identifies its operator with
\(g_{\mathrm{YM}}^{-2}\mathrm d_1^*\mathrm d_1\) on precisely the domain
displayed in \cref{cusp-gauge-continuum-hessian-operator}.
\end{proof}

\subsection{The finite cubical cochain complex}
\label{cusp-gauge-finite-subsection}

Fix \(N\ge2\) and put
\[
 \Lambda_N=(\mathbb Z/N\mathbb Z)^4.
\]
The vertex \(n\in\Lambda_N\) is realized in \(F_T\) as the class of
\(P_T(n/N)\).  Let \(E=(\mathbb R^4)^*\), with its Euclidean inner product,
and let \(\mathbf e_1,\ldots,\mathbf e_4\) be the standard coordinate
columns in the lattice-coordinate space, also used as the coefficient
columns of its dual coordinate basis.  Define the physical coframe
\begin{equation}
 \vartheta_T^j=P_T^{-\mathsf T}\mathbf e_j\in E,
 \qquad 1\le j\le4.
 \label{cusp-gauge-finite-physical-coframe}
\end{equation}
Thus \(\vartheta_T^j=\mathrm dy^j\) in physical \(x\)-coordinates, and the
Gram matrix of this coframe is \(G_T^{-1}\).

For \(0\le p\le4\), define
\begin{equation}
 C_N^p(T;\mathfrak g_{\mathbb C})
 =\operatorname{Map}\left(
   \Lambda_N,\bigwedge^pE_{\mathbb C}\otimes\mathfrak g_{\mathbb C}
  \right).
 \label{cusp-gauge-finite-cochain-spaces}
\end{equation}
These spaces are finite-dimensional.  Every finite differential, adjoint,
Hodge operator, projection, and reduced Hessian below therefore acts on its
entire displayed finite-dimensional source space; no additional operator
domain is implicit.
The coefficient of
\(\vartheta_T^{i_1}\wedge\cdots\wedge\vartheta_T^{i_p}\) at \(n\) is the
coefficient assigned to the positively oriented coordinate \(p\)-cell
based at \(n\).  Let
\[
 (\mathsf T_ja)(n)=a(n+\mathbf e_j)
\]
and define the forward and backward differences by
\begin{equation}
 D_{N,j}^+=N(\mathsf T_j-\operatorname{id}),\qquad
 (D_{N,j}^-a)(n)
 =N\bigl(a(n)-a(n-\mathbf e_j)\bigr).
 \label{cusp-gauge-finite-differences}
\end{equation}
The finite exterior derivative is
\begin{equation}
 \mathrm d_Na
 =\sum_{j=1}^4
   \varepsilon(\vartheta_T^j)D_{N,j}^+a.
 \label{cusp-gauge-finite-exterior-derivative}
\end{equation}

\begin{lemma}[Cochain property and induced exterior metric]
\label{cusp-gauge-finite-complex-metric}
The maps in
\begin{equation}
 0\longrightarrow C_N^0
 \mathrel{\mathop{\longrightarrow}^{\mathrm d_N}} C_N^1
 \mathrel{\mathop{\longrightarrow}^{\mathrm d_N}} C_N^2
 \mathrel{\mathop{\longrightarrow}^{\mathrm d_N}} C_N^3
 \mathrel{\mathop{\longrightarrow}^{\mathrm d_N}} C_N^4
 \longrightarrow0
 \label{cusp-gauge-finite-cochain-complex}
\end{equation}
form a cochain complex.  Equip \(C_N^p\) with
\begin{equation}
 \langle a,b\rangle_{N,p,T}
 =\frac{D_T}{N^4}\sum_{n\in\Lambda_N}
 \langle a(n),b(n)\rangle_{\wedge^pE\otimes\mathfrak g_{\mathbb C}}.
 \label{cusp-gauge-finite-inner-product}
\end{equation}
For increasing multi-indices
\(I=(i_1,\ldots,i_p)\) and \(J=(j_1,\ldots,j_p)\), the exterior metric is
exactly
\begin{equation}
 \left\langle
  \vartheta_T^{i_1}\wedge\cdots\wedge\vartheta_T^{i_p},
  \vartheta_T^{j_1}\wedge\cdots\wedge\vartheta_T^{j_p}
 \right\rangle
 =\det\left((G_T^{-1})_{i_a j_b}\right)_{a,b=1}^p.
 \label{cusp-gauge-finite-exterior-metric}
\end{equation}
\end{lemma}

\begin{proof}
The differences \(D_{N,j}^+\) commute, whereas the exterior
multiplications \(\varepsilon(\vartheta_T^j)\) anticommute.  Therefore
\[
 \mathrm d_N^2
 =\sum_{j,\ell}
 \varepsilon(\vartheta_T^j)\varepsilon(\vartheta_T^\ell)
 D_{N,j}^+D_{N,\ell}^+=0,
\]
which proves that \cref{cusp-gauge-finite-cochain-complex} is a complex.
The factor \(D_T/N^4\) in
\cref{cusp-gauge-finite-inner-product} is the exact Riemann-sum weight:
\(D_T\) is the covolume of \(P_T\mathbb Z^4\).  The Euclidean inner
product of \(\vartheta_T^i=P_T^{-\mathsf T}\mathbf e_i\) and
\(\vartheta_T^j=P_T^{-\mathsf T}\mathbf e_j\) is
\((G_T^{-1})_{ij}\).  The induced inner product on an exterior power is
the determinant of the matrix of pairings, proving
\cref{cusp-gauge-finite-exterior-metric}.
\end{proof}

\begin{remark}[Integrated cellular cochains give the same operator]
\label{cusp-gauge-finite-integrated-normalization}
The convention above stores form coefficients.  If \(c_I(n)\) instead
stores the integral over a \(p\)-cell, put \(a_I(n)=N^pc_I(n)\).  The
unscaled cellular coboundary \(\mathsf T_j-\operatorname{id}\), together
with the inner product
\[
 \langle c,c'\rangle^{\mathrm{int}}_{N,p,T}
 =D_TN^{2p-4}\sum_{n\in\Lambda_N}
 \langle c(n),c'(n)\rangle_{\wedge^pE\otimes\mathfrak g_{\mathbb C}},
\]
is carried isometrically to
\cref{cusp-gauge-finite-exterior-derivative,cusp-gauge-finite-inner-product}.
Thus the two conventions give identical Hessian symbols and eigenvalues.
\end{remark}

\begin{proposition}[Adjoint and Fourier wedge--contraction symbols]
\label{cusp-gauge-finite-symbols}
The adjoint of \(\mathrm d_N:C_N^p\to C_N^{p+1}\) is
\begin{equation}
 \mathrm d_N^*
 =-\sum_{j=1}^4
 \iota_{(\vartheta_T^j)^\sharp}D_{N,j}^-.
 \label{cusp-gauge-finite-adjoint}
\end{equation}
For \(k\in\Lambda_N\), define
\begin{align}
 \chi_{N,k}(n)&=\exp(2\pi i k\cdot n/N),\notag\\
 \mathfrak d_N(k)_j&=N\left(e^{2\pi i k_j/N}-1\right),
 \label{cusp-gauge-finite-coordinate-symbol}\\
 \zeta_{N,T}(k)&=\sum_{j=1}^4\mathfrak d_N(k)_j\vartheta_T^j
 =P_T^{-\mathsf T}\mathfrak d_N(k)\in E_{\mathbb C}.
 \label{cusp-gauge-finite-physical-symbol}
\end{align}
For every constant exterior amplitude \(\alpha\),
\begin{align}
 \mathrm d_N(\chi_{N,k}\alpha)
 &=\chi_{N,k}\varepsilon(\zeta_{N,T}(k))\alpha,
 \label{cusp-gauge-finite-d-symbol}\\
 \mathrm d_N^*(\chi_{N,k}\alpha)
 &=\chi_{N,k}
   \iota_{\overline{\zeta_{N,T}(k)}^\sharp}\alpha,
 \label{cusp-gauge-finite-dstar-symbol}\\
 (\mathrm d_N\mathrm d_N^*+\mathrm d_N^*\mathrm d_N)
 (\chi_{N,k}\alpha)
 &=\lvert\zeta_{N,T}(k)\rvert^2\chi_{N,k}\alpha.
 \label{cusp-gauge-finite-hodge-symbol}
\end{align}
Consequently the harmonic \(p\)-cochains are exactly the constant
cochains and have complex dimension \(\binom4p d_G\).
\end{proposition}

\begin{proof}
Translation is unitary for
\cref{cusp-gauge-finite-inner-product}; hence
\[
 (D_{N,j}^+)^*
 =N(\mathsf T_j^{-1}-\operatorname{id})=-D_{N,j}^-.
\]
For the induced exterior metric, the adjoint of exterior multiplication
by \(\vartheta_T^j\) is contraction by
\((\vartheta_T^j)^\sharp\).  This proves
\cref{cusp-gauge-finite-adjoint}.  On \(\chi_{N,k}\), the forward
difference has eigenvalue \(\mathfrak d_N(k)_j\), whereas \(-D_{N,j}^-\) has
eigenvalue \(\overline{\mathfrak d_N(k)_j}\).  Summing proves
\cref{cusp-gauge-finite-d-symbol,cusp-gauge-finite-dstar-symbol}.  The
Hermitian exterior-algebra identity
\[
 \iota_{\overline\zeta^\sharp}\varepsilon(\zeta)
 +\varepsilon(\zeta)\iota_{\overline\zeta^\sharp}
 =\lvert\zeta\rvert^2\operatorname{id}
\]
proves \cref{cusp-gauge-finite-hodge-symbol}.  Finally, \(\mathfrak d_N(k)=0\)
exactly when \(k=0\) in \(\Lambda_N\), and \(P_T\) is invertible.  Thus the
Hodge symbol vanishes exactly on the zero Fourier mode, where the entire
constant exterior amplitude is harmonic.
\end{proof}

Define the finite reduced Coulomb slice by
\begin{equation}
 \mathscr C_{N,T}
 =\ker(\mathrm d_N^*:C_N^1\to C_N^0)
  \cap(\ker\Delta_{N,T}^{(1)})^\perp,
 \qquad
 \Delta_{N,T}^{(1)}
 =\mathrm d_N\mathrm d_N^*+\mathrm d_N^*\mathrm d_N.
 \label{cusp-gauge-finite-reduced-coulomb}
\end{equation}
Equivalently, one deletes the zero Fourier mode and, for each \(k\ne0\),
restricts the amplitude to
\begin{equation}
 \mathcal T_{N,T,k}
 =\ker\iota_{\overline{\zeta_{N,T}(k)}^\sharp}
 \subset E_{\mathbb C}.
 \label{cusp-gauge-finite-transverse-space}
\end{equation}
The exact orthogonal projection is
\begin{equation}
 \Pi_{N,T,k}\alpha
 =\alpha-\zeta_{N,T}(k)
 \frac{
  \iota_{\overline{\zeta_{N,T}(k)}^\sharp}\alpha
 }{\lvert\zeta_{N,T}(k)\rvert^2}.
 \label{cusp-gauge-finite-coulomb-projection}
\end{equation}
Thus the projection uses the physical covector
\(P_T^{-\mathsf T}\mathfrak d_N(k)\), including every skew-metric cross term.

The finite quadratic gauge form and its associated operator are
\begin{equation}
 Q_{N,T}(a)
 =\frac1{2g_{\mathrm{YM}}^2}
   \lVert\mathrm d_Na\rVert_{N,2,T}^2,
 \qquad
 \mathcal L_{N,T}
 =g_{\mathrm{YM}}^{-2}\mathrm d_N^*\mathrm d_N.
 \label{cusp-gauge-finite-quadratic-form}
\end{equation}
This is the exact linearized, Lie-algebra-valued cubical discretization of
\cref{cusp-gauge-continuum-hessian-form}; it is not a nonlinear
compact-link Wilson action.

\begin{corollary}[Exact finite Coulomb spectrum]
\label{cusp-gauge-finite-spectrum}
On the \(k\ne0\) Fourier fibre of \(\mathscr C_{N,T}\),
\(\mathcal L_{N,T}\) has the single eigenvalue
\begin{equation}
 \lambda_{N,T}^{\mathrm{YM}}(k)
 =\frac1{g_{\mathrm{YM}}^2}
  \lvert P_T^{-\mathsf T}\mathfrak d_N(k)\rvert^2
 =\frac1{g_{\mathrm{YM}}^2}
  \mathfrak d_N(k)^*G_T^{-1}\mathfrak d_N(k)
 \label{cusp-gauge-finite-mode-eigenvalue}
\end{equation}
with complex multiplicity \(3d_G\).
\end{corollary}

\begin{proof}
On \(\mathcal T_{N,T,k}\), the contraction term in
\[
 \iota_{\overline\zeta^\sharp}\varepsilon(\zeta)\alpha
 =\lvert\zeta\rvert^2\alpha
  -\varepsilon(\zeta)\iota_{\overline\zeta^\sharp}\alpha
\]
vanishes.  This proves the eigenvalue.  Since \(\zeta_{N,T}(k)\ne0\), the
contraction has rank one on \(E_{\mathbb C}\), and its kernel has dimension
three.  Tensoring with \(\mathfrak g_{\mathbb C}\) gives the stated
multiplicity.
\end{proof}

\subsection{Exact coordinate-mode eigenvalues and all mesh regimes}
\label{cusp-gauge-finite-limits-subsection}

Set
\begin{equation}
 \rho_N=4N^2\sin^2\left(\frac\pi N\right).
 \label{cusp-gauge-finite-rho}
\end{equation}
Because \(\mathfrak d_N(\mathbf e_\gamma)\) and \(\mathfrak d_N(\mathbf e_u)\) are the common
scalar \(N(e^{2\pi i/N}-1)\) times the corresponding coordinate vector,
\cref{cusp-gauge-finite-mode-eigenvalue,cusp-gauge-short-mode-lengths}
give the exact eigenvalues
\begin{align}
 \lambda_{N,T}^{\mathrm{YM}}(\mathbf e_\gamma)
 &=\frac{\rho_N}{g_{\mathrm{YM}}^2}
   \frac{L_T^2+m_T^2}{D_T^2},
 \label{cusp-gauge-finite-gamma-eigenvalue}\\
 \lambda_{N,T}^{\mathrm{YM}}(\mathbf e_u)
 &=\frac{\rho_N}{g_{\mathrm{YM}}^2}
   \frac{q_T^2+36m_T^2}{D_T^2}.
 \label{cusp-gauge-finite-u-eigenvalue}
\end{align}
Each individual complex Fourier fibre contributes multiplicity \(3d_G\).
For real cochains, the pair
\(\{\mathbf e_j,-\mathbf e_j\}\) contributes \(6d_G\) dimensions when
\(N>2\).  When \(N=2\), these two labels coincide, so its real contribution
is \(3d_G\).

\begin{theorem}[Finite gap and simultaneous cusp--mesh limits]
\label{cusp-gauge-finite-gap-limits}
There is \(T_1\), independent of \(N\ge2\), such that for \(T\ge T_1\) the
first eigenvalue of \(\mathcal L_{N,T}\) on \(\mathscr C_{N,T}\) is
\begin{equation}
 \gamma_{N,T}
 =\min\left\{
   \lambda_{N,T}^{\mathrm{YM}}(\mathbf e_\gamma),
   \lambda_{N,T}^{\mathrm{YM}}(\mathbf e_u)
 \right\}.
 \label{cusp-gauge-finite-gap-exact}
\end{equation}
Uniformly for \(N\ge2\),
\begin{equation}
 \gamma_{N,T}
 =\frac{\rho_N}{g_{\mathrm{YM}}^2T^2}
 \left(
  1+\frac{-2H+2\min\{B,0\}}T+O(T^{-2})
 \right).
 \label{cusp-gauge-finite-gap-expansion}
\end{equation}
Let
\[
 I_{N,T}=\{j\in\{\gamma,u\}:
 \lambda_{N,T}^{\mathrm{YM}}(\mathbf e_j)=\gamma_{N,T}\},
 \qquad
 \eta_N=\begin{cases}1,&N=2,\\2,&N>2.\end{cases}
\]
The real multiplicity is exactly
\begin{equation}
 3d_G\eta_N\,\#I_{N,T}.
 \label{cusp-gauge-finite-gap-multiplicity}
\end{equation}
The limiting conclusions are:
\begin{enumerate}[label=(\roman*)]
\item For fixed \(N\ge2\),
\begin{equation}
 T^2\gamma_{N,T}
 \longrightarrow\frac{\rho_N}{g_{\mathrm{YM}}^2},
 \qquad \gamma_{N,T}\longrightarrow0.
 \label{cusp-gauge-finite-fixed-N}
\end{equation}
\item For every integer-valued \(N(T)\ge2\), with no growth assumption,
there are constants \(0<c<C<\infty\), independent of \(N(T)\), such that
\begin{equation}
 \frac{c}{T^2}\le\gamma_{N(T),T}\le\frac{C}{T^2}
 \label{cusp-gauge-finite-arbitrary-N}
\end{equation}
for all large \(T\).  Hence \(\gamma_{N(T),T}\to0\), although the scaled
gaps need not converge when \(N(T)\) oscillates.
\item If \(N(T)\to\infty\), then
\begin{equation}
 \gamma_{N(T),T}
 =\frac{4\pi^2}{g_{\mathrm{YM}}^2T^2}
  \left(1+O(T^{-1})+O(N(T)^{-2})\right).
 \label{cusp-gauge-finite-growing-N}
\end{equation}
\item The largest physical elementary edge has length
\begin{equation}
 \max_{1\le j\le4}
 \frac{\lVert P_T\mathbf e_j\rVert}{N(T)}
 =\frac{T}{N(T)}(1+O(T^{-1})).
 \label{cusp-gauge-finite-mesh-size}
\end{equation}
Thus the varying continuum torus is resolved if and only if
\(N(T)/T\to\infty\).  This refinement condition is not needed for the
algebraic soft-mode conclusion in (ii).
\end{enumerate}
\end{theorem}

\begin{proof}
Let
\[
 \sigma_N=2N\sin(\pi/N),\qquad \rho_N=\sigma_N^2.
\]
The least nonzero modulus of a coordinate of \(\mathfrak d_N(k)\) is \(\sigma_N\),
attained precisely at frequency \(\pm1\), with the two signs equal for
\(N=2\).  If \(\mathfrak d_N(k)\) is supported in its first two components but
\(k\notin\{\pm\mathbf e_\gamma,\pm\mathbf e_u\}\), then
\begin{equation}
 |\mathfrak d_N(k)_1|^2+|\mathfrak d_N(k)_2|^2\ge2\sigma_N^2.
 \label{cusp-gauge-finite-next-frequency}
\end{equation}
Indeed, either both coordinates are nonzero, or the single nonzero
frequency is not \(\pm1\).  In the latter case
\[
 \frac{\sin(2\pi/N)}{\sin(\pi/N)}
 =2\cos(\pi/N)\ge\sqrt2
\]
for \(N\ge4\), and the cases \(N=2,3\) are immediate.

For \(z=(z_1,z_2,0,0)^{\mathsf t}\in\mathbb C^4\), the exact inverse
transpose in \cref{cusp-spec-exact-dual-lattice} gives
\begin{equation}
 P_T^{-\mathsf T}z
 =\left(
 0,\frac{m_Tz_1+q_Tz_2}{D_T},
 0,\frac{-L_Tz_1+6m_Tz_2}{D_T}
 \right).
 \label{cusp-gauge-finite-soft-block}
\end{equation}
The scaled \(2\times2\) matrix
\[
 \frac{T}{D_T}
 \begin{pmatrix}m_T&q_T\\-L_T&6m_T\end{pmatrix}
 =
 \begin{pmatrix}0&1\\-1&0\end{pmatrix}
 +O_{\mathrm{op}}(T^{-1})
\]
has all singular values \(1+O(T^{-1})\), also on complex vectors.  Thus
the factor \(\sqrt2\) in
\cref{cusp-gauge-finite-next-frequency} separates every other soft-plane
frequency from the coordinate frequencies uniformly in \(N\) for large
\(T\).

For a general \(z\in\mathbb C^4\), the full inverse transpose in
\cref{cusp-spec-exact-dual-lattice} has first and third components
\(z_3,z_4\).  Hence a finite frequency with nonzero third or fourth
coordinate satisfies
\[
 \lVert P_T^{-\mathsf T}\mathfrak d_N(k)\rVert^2
 \ge |\mathfrak d_N(k)_3|^2+|\mathfrak d_N(k)_4|^2
 \ge\sigma_N^2,
\]
whereas either coordinate candidate is
\(O(\sigma_N^2/T^2)\).  This proves the exact minimizing list and
\cref{cusp-gauge-finite-gap-exact}; the reality discussion before the
theorem proves \cref{cusp-gauge-finite-gap-multiplicity}.

Multiplying the two expansions in
\cref{cusp-spec-shortest-modes} by
\(\rho_N/g_{\mathrm{YM}}^2\) proves
\cref{cusp-gauge-finite-gap-expansion}.  The elementary sine bounds
\(2x/\pi\le\sin x\le x\) for \(0<x\le\pi/2\) give
\begin{equation}
 16\le\rho_N\le4\pi^2.
 \label{cusp-gauge-finite-rho-bounds}
\end{equation}
Together with \cref{cusp-gauge-finite-gap-expansion}, these bounds prove
(i) and (ii).  The Taylor expansion
\(\rho_N=4\pi^2(1+O(N^{-2}))\) proves (iii).

The first two columns of \(P_T\), displayed in
\cref{cusp-spec-exact-dual-lattice}, have Euclidean lengths
\(T(1+O(T^{-1}))\), while its last two columns have length one.  Division
by \(N(T)\) proves \cref{cusp-gauge-finite-mesh-size} and (iv).
\end{proof}

\begin{remark}[Every fixed reduced model remains positively gapped]
\label{cusp-gauge-pointwise-versus-uniform}
For each fixed \(T\), the reduced continuum operator has a positive first
eigenvalue.  For each fixed pair \((T,N)\), the reduced finite operator
also has a positive first eigenvalue.  What fails is a positive lower
bound uniform as \(T\to\infty\).  Before reduction, gauge-longitudinal
directions and constant harmonic one-forms are separate Hessian-zero
directions.
\end{remark}

\subsection{Continuum Coulomb reduction and exact spectrum}
\label{cusp-gauge-continuum-spectrum-subsection}

The infinitesimal gauge directions
\(\mathrm d\Omega^0(F_T;\mathfrak g)\) and the harmonic one-forms are
Hessian-zero directions.  Coulomb gauge removes the former but not the
latter.  The reduced Coulomb space is therefore
\begin{equation}
 \mathscr C_T=
 \left\{
 a\in L^2\Omega^1(F_T;\mathfrak g):
 \mathrm d^*a=0\ \text{distributionally},\
 a\perp\mathcal H^1(F_T;\mathfrak g)
 \right\}.
 \label{cusp-gauge-continuum-reduced-coulomb}
\end{equation}
The subspace \(\mathscr C_T\) reduces \(\mathcal L_T\).  On it the closed
Hessian form has domain
\(H^1\Omega^1(F_T;\mathfrak g)\cap\mathscr C_T\), and the restricted
self-adjoint operator has domain
\(H^2\Omega^1(F_T;\mathfrak g)\cap\mathscr C_T\).  Indeed, on the
distributionally coclosed slice the graph norm of \(\mathrm d_1\), together
with the \(L^2\)-norm, is the Hodge graph norm
\[
 \lVert a\rVert_{L^2}^2+
 \lVert\mathrm da\rVert_{L^2}^2+
 \lVert\mathrm d^*a\rVert_{L^2}^2.
\]
The Gaffney estimate on the compact flat torus identifies its completion
with \(H^1\), and elliptic regularity for the Hodge Laplacian identifies the
operator domain with \(H^2\).  Orthogonality to the finite-dimensional
harmonic space is preserved by the Fourier multiplier.

For \(k\in\mathbb Z^4\), write
\[
 \xi_{T,k}=P_T^{-\mathsf T}k,\qquad
 \chi_{T,k}([x])
 =\exp\bigl(2\pi i\langle\xi_{T,k},x\rangle\bigr).
\]
For a covector \(\xi\), let \(\varepsilon(\xi)\) denote exterior
multiplication by \(\xi\), and use the Euclidean metric to define
\(\xi^\sharp\) and contraction \(\iota_{\xi^\sharp}\).

\begin{proposition}[Continuum Fourier symbol, slice, and multiplicity]
\label{cusp-gauge-continuum-symbol}
For
\(\alpha\in(\mathbb R^4)^*_{\mathbb C}\otimes\mathfrak g_{\mathbb C}\),
\begin{align}
 \mathrm d(\chi_{T,k}\alpha)
 &=2\pi i\,\chi_{T,k}\varepsilon(\xi_{T,k})\alpha,
 \label{cusp-gauge-continuum-d-symbol}\\
 \mathrm d^*(\chi_{T,k}\alpha)
 &=-2\pi i\,\chi_{T,k}\iota_{\xi_{T,k}^\sharp}\alpha,
 \label{cusp-gauge-continuum-dstar-symbol}\\
 \mathrm d^*\mathrm d(\chi_{T,k}\alpha)
 &=4\pi^2\chi_{T,k}
 \left(
  \lvert\xi_{T,k}\rvert^2\alpha
  -\varepsilon(\xi_{T,k})
   \iota_{\xi_{T,k}^\sharp}\alpha
 \right).
 \label{cusp-gauge-continuum-dstard-symbol}
\end{align}
For \(k\ne0\), the Coulomb polarization space is
\begin{equation}
 \mathcal T_{T,k}
 =\ker\iota_{\xi_{T,k}^\sharp}
 \subset(\mathbb R^4)^*_{\mathbb C}.
 \label{cusp-gauge-continuum-transverse-space}
\end{equation}
It has complex dimension three.  On
\(\mathcal T_{T,k}\otimes\mathfrak g_{\mathbb C}\), the Hessian
eigenvalue is
\begin{equation}
 \lambda_T^{\mathrm{YM}}(k)
 =\frac{4\pi^2}{g_{\mathrm{YM}}^2}
 \lvert P_T^{-\mathsf T}k\rvert^2
 \label{cusp-gauge-continuum-mode-eigenvalue}
\end{equation}
with complex multiplicity \(3d_G\) on that signed Fourier fibre.  If
\[
 K_T(\lambda)
 =\{k\in\mathbb Z^4\setminus\{0\}:
       \lambda_T^{\mathrm{YM}}(k)=\lambda\},
\]
the total complex multiplicity, and also the dimension of the real
eigenspace, is
\begin{equation}
 3d_G\,\#K_T(\lambda).
 \label{cusp-gauge-continuum-shell-multiplicity}
\end{equation}
Equivalently, a conjugate pair \(\{k,-k\}\) contributes \(6d_G\) real
dimensions.
\end{proposition}

\begin{proof}
Differentiating the character proves
\cref{cusp-gauge-continuum-d-symbol}.  On the Euclidean cover,
\[
 \mathrm d^*
 =-\sum_{j=1}^4
 \iota_{\partial_{x_j}}\partial_{x_j},
\]
which proves \cref{cusp-gauge-continuum-dstar-symbol}.  The exterior
algebra identity
\[
 \iota_{\xi^\sharp}\varepsilon(\xi)
 +\varepsilon(\xi)\iota_{\xi^\sharp}
 =\lvert\xi\rvert^2\operatorname{id}
\]
proves \cref{cusp-gauge-continuum-dstard-symbol}.
For \(k\ne0\), the invertibility of \(P_T\) gives \(\xi_{T,k}\ne0\).
Contraction by \(\xi_{T,k}^\sharp\) has rank one on one-forms, so its
kernel has dimension three.  On that kernel the second term of
\cref{cusp-gauge-continuum-dstard-symbol} vanishes, proving the eigenvalue
and its per-fibre multiplicity.  The zero Fourier mode consists of the
constant harmonic one-forms and has already been removed in
\cref{cusp-gauge-continuum-reduced-coulomb}.  Fourier completeness gives
all eigenvectors.  Complex conjugation exchanges \(k\) with \(-k\);
the reality condition \(\alpha_{-k}=\overline{\alpha_k}\) turns each
conjugate pair of complex \(3d_G\)-dimensional spaces into a real
\(6d_G\)-dimensional space.  This proves
\cref{cusp-gauge-continuum-shell-multiplicity}.
\end{proof}

\begin{remark}[The Coulomb projection uses the physical covector metric]
\label{cusp-gauge-continuum-projection}
In lattice-coordinate one-form coefficients \(a_y\), the mode \(k\) obeys
\begin{equation}
 k^{\mathsf T}G_T^{-1}a_y=0,\qquad
 \Pi_{T,k}a_y
 =a_y-k\frac{k^{\mathsf T}G_T^{-1}a_y}
 {k^{\mathsf T}G_T^{-1}k}.
 \label{cusp-gauge-continuum-coordinate-projection}
\end{equation}
After conversion to physical \(x\)-coefficients, this is Euclidean
projection perpendicular to \(P_T^{-\mathsf T}k\).  It is generally
incorrect to impose Euclidean coefficient orthogonality to \(k\), or to
delete one coordinate merely because \(k\) is a coordinate lattice vector.
\end{remark}

\begin{theorem}[Exact reduced continuum Hessian gap]
\label{cusp-gauge-continuum-gap}
For all sufficiently large \(T\), the first eigenvalue of
\(\mathcal L_T\) on \(\mathscr C_T\) is
\begin{equation}
 \gamma_T
 =\frac{4\pi^2}{g_{\mathrm{YM}}^2}
 \min\left\{
   \frac{L_T^2+m_T^2}{D_T^2},
   \frac{q_T^2+36m_T^2}{D_T^2}
 \right\}.
 \label{cusp-gauge-continuum-gap-exact}
\end{equation}
It has the radial expansion
\begin{equation}
 \gamma_T
 =\frac{4\pi^2}{g_{\mathrm{YM}}^2T^2}
 \left(
  1+\frac{-2H+2\min\{B,0\}}T+O(T^{-2})
 \right).
 \label{cusp-gauge-continuum-gap-expansion}
\end{equation}
Let
\[
 I_T=\{j\in\{\gamma,u\}:a_{j,T}
             =\min\{a_{\gamma,T},a_{u,T}\}\}.
\]
The real multiplicity of \(\gamma_T\) is exactly
\begin{equation}
 6d_G\,\#I_T.
 \label{cusp-gauge-continuum-gap-multiplicity}
\end{equation}
\end{theorem}

\begin{proof}
By \cref{cusp-gauge-continuum-symbol}, minimizing the reduced Hessian
eigenvalue is exactly the shortest-nonzero-dual-vector problem.  The
shortest-mode theorem
\cref{cusp-spec-shortest-modes} says that, for all large \(T\), the only
possible minimizers are the two signed pairs with exact squared lengths
given in \cref{cusp-gauge-short-mode-lengths}.  This proves
\cref{cusp-gauge-continuum-gap-exact}.  The expansions already proved in
\cref{cusp-spec-shortest-modes} give relative corrections
\(2(B-H)\) and \(-2H\); their minimum is
\(-2H+2\min\{B,0\}\), proving
\cref{cusp-gauge-continuum-gap-expansion}.  Each minimizing signed pair
contributes \(6d_G\) real dimensions by
\cref{cusp-gauge-continuum-symbol}, which proves the multiplicity formula.
\end{proof}

\subsection{The physical boundary and the exact no-gap bridge}
\label{cusp-gauge-physical-boundary-subsection}

The operators proved above are a classical Hessian and a finite quadratic
cochain approximation on the internal flat manifolds \(F_T\).  They do not
construct a quantum gauge theory on standard spacetime \(\mathbb R^4\), and
they do not determine the mass spectrum of such a theory.  The precise
boundaries are as follows.
\begin{enumerate}[label=(\arabic*)]
\item The unscaled cusp retains two compact geometric directions.  The
period parameter \(T\) has not been identified with a spacetime box size,
a lattice spacing, a correlation length, or a renormalization scale.
\item The operator \(g_{\mathrm{YM}}^{-2}\mathrm d^*\mathrm d\) is the
quadratic classical fluctuation operator at the product connection.  It
contains no nonperturbative interacting measure, dimensional
transmutation, or confinement argument.
\item The spaces \(C_N^p(T;\mathfrak g_{\mathbb C})\) are linear
Lie-algebra cochains.  They are not compact-group link variables, and
\cref{cusp-gauge-finite-quadratic-form} is not a Wilson plaquette action.
\item No Euclidean time direction or reflection, reflection-positive
measure, transfer matrix, Osterwalder--Schrader reconstruction, Gauss-law
physical quotient, local observable algebra, thermodynamic limit, or
continuum limit in fixed physical units has been constructed from these
operators.
\item An eigenvalue of a four-dimensional elliptic Hessian has units of
squared inverse length.  Without a time reconstruction and a scale
conversion, it is not an eigenvalue of a physical Hamiltonian.  In
particular, a dimensionless transfer-matrix gap can behave as
\(a_tm\to0\) while the physical mass \(m>0\) remains fixed.
\end{enumerate}

The exact extra statement that would transfer the auxiliary soft sequence
to one already constructed physical Hamiltonian is the following.  It
must use the full zero-energy sector rather than only a chosen vacuum
vector.

\begin{proposition}[Typed state-and-energy no-gap bridge]
\label{cusp-gauge-no-gap-bridge}
Let \(\mathcal H_{\mathrm{phys}}\) be a Hilbert space and let \(H\ge0\) be
self-adjoint on it.  Its closed energy form is
\begin{equation}
 \mathfrak h[\psi]=\lVert H^{1/2}\psi\rVert^2,\qquad
 \mathcal D(\mathfrak h)=\mathcal D(H^{1/2}).
 \label{cusp-gauge-physical-energy-form}
\end{equation}
Let
\begin{equation}
 P_0=\mathbf 1_{\{0\}}(H),\qquad
 \mathcal V=\ker H=P_0\mathcal H_{\mathrm{phys}},\qquad
 \mathcal H_{\mathrm{exc}}
 =(\operatorname{id}-P_0)\mathcal H_{\mathrm{phys}},
 \label{cusp-gauge-vacuum-excitation-spaces}
\end{equation}
and fix a normalized vacuum vector \(\Omega\in\mathcal V\).  Then
\begin{equation}
 \mathcal H_{\mathrm{phys}}
 =\mathcal V\oplus\mathcal H_{\mathrm{exc}}
 \label{cusp-gauge-reducing-decomposition}
\end{equation}
reduces \(H\).  Denote its excitation restriction by
\[
 H_{\mathrm{exc}}
 =H|_{\mathcal D(H)\cap\mathcal H_{\mathrm{exc}}}.
\]

For every \(n\), let \(\mathcal K_n\) be an auxiliary Hilbert space with a
closed nonnegative quadratic form \(\mathfrak q_n\).  Let
\(u_n\in\mathcal D(\mathfrak q_n)\) satisfy
\begin{equation}
 \lVert u_n\rVert_{\mathcal K_n}=1,\qquad
 \mathfrak q_n[u_n]\longrightarrow0.
 \label{cusp-gauge-auxiliary-soft-states}
\end{equation}
Suppose there are linear maps
\begin{equation}
 J_n:\mathcal D(\mathfrak q_n)\longrightarrow\mathcal D(\mathfrak h),
 \qquad
 \psi_n=(\operatorname{id}-P_0)J_nu_n,
 \label{cusp-gauge-bridge-maps}
\end{equation}
and numbers \(c>0\), \(C_n\ge0\), and \(\varepsilon_n\ge0\) such that
\begin{align}
 \lVert\psi_n\rVert&\ge c,
 \label{cusp-gauge-bridge-norm}\\
 \mathfrak h[\psi_n]
 &\le C_n\mathfrak q_n[u_n]+\varepsilon_n,
 \qquad
 C_n\mathfrak q_n[u_n]+\varepsilon_n\longrightarrow0.
 \label{cusp-gauge-bridge-energy}
\end{align}
Then
\begin{equation}
 \inf\operatorname{Spec}(H_{\mathrm{exc}})=0
 \label{cusp-gauge-bridge-zero-infimum}
\end{equation}
and, for every \(\epsilon>0\),
\begin{equation}
 \mathbf 1_{(0,\epsilon)}(H)\ne0.
 \label{cusp-gauge-bridge-spectrum-near-zero}
\end{equation}
Moreover,
\begin{equation}
 0\in\operatorname{Spec}(H_{\mathrm{exc}}),\qquad
 \ker H_{\mathrm{exc}}=\{0\}.
 \label{cusp-gauge-bridge-continuous-threshold}
\end{equation}
Thus zero is a continuous spectral threshold, essential under the usual
Weyl convention, and \(H\) has no positive gap above its full zero-energy
sector.

If the vacuum is unique, so that \(\mathcal V=\mathbb C\Omega\), one may
replace the projection in \cref{cusp-gauge-bridge-maps} by the requirement
\(J_nu_n\perp\Omega\).  Without uniqueness, orthogonality only to
\(\Omega\) is insufficient.
\end{proposition}

\begin{proof}
The operator \(P_0\) is a spectral projection of \(H\), so it commutes
with the resolvent and with the spectral calculus of \(H\).  Hence
\cref{cusp-gauge-reducing-decomposition} is a reducing decomposition and
\[
 \mathcal D(\mathfrak h|_{\mathcal H_{\mathrm{exc}}})
 =\mathcal D(H^{1/2})\cap\mathcal H_{\mathrm{exc}}.
\]
The normalized vectors
\[
 v_n=\frac{\psi_n}{\lVert\psi_n\rVert}
\]
belong to this excitation form domain.  By
\cref{cusp-gauge-bridge-norm,cusp-gauge-bridge-energy},
\[
 \mathfrak h[v_n]
 \le c^{-2}
 \bigl(C_n\mathfrak q_n[u_n]+\varepsilon_n\bigr)
 \longrightarrow0.
\]
The form version of the spectral variational identity is
\[
 \inf\operatorname{Spec}(H_{\mathrm{exc}})
 =\inf_{\substack{
 0\ne v\in\mathcal D(H^{1/2})\cap\mathcal H_{\mathrm{exc}}
 }}
 \frac{\mathfrak h[v]}{\lVert v\rVert^2}.
\]
Nonnegativity and the sequence \(v_n\) prove
\cref{cusp-gauge-bridge-zero-infimum}.

If \(\mathbf 1_{(0,\epsilon)}(H)=0\) for some \(\epsilon>0\), the spectral
measure of every excitation vector would be supported in
\([\epsilon,\infty)\), and therefore
\[
 \mathfrak h[v]\ge\epsilon\lVert v\rVert^2
\]
throughout the excitation form domain.  This contradicts
\(\mathfrak h[v_n]\to0\), proving
\cref{cusp-gauge-bridge-spectrum-near-zero}.  Since the spectrum of a
self-adjoint operator is closed,
\cref{cusp-gauge-bridge-zero-infimum} puts zero in
\(\operatorname{Spec}(H_{\mathrm{exc}})\).  Its kernel is trivial by the
definition of \(\mathcal H_{\mathrm{exc}}\).  A self-adjoint operator has
no residual spectrum, and
\cref{cusp-gauge-bridge-spectrum-near-zero} makes zero an accumulation
point.  This proves
\cref{cusp-gauge-bridge-continuous-threshold} and the essential-threshold
assertion.

If \(\mathcal V=\mathbb C\Omega\), orthogonality to \(\Omega\) is exactly
membership in \(\mathcal H_{\mathrm{exc}}\).  If the kernel is larger, a
zero-energy vector can be orthogonal to the selected \(\Omega\), so that
orthogonality alone does not produce an excitation vector.  This proves
the final assertion.
\end{proof}

\begin{remark}[The present geometry supplies only the auxiliary side]
\label{cusp-gauge-missing-physical-bridge}
Choose any \(T_n\to\infty\), any \(N_n\ge2\), and a unit real transverse
eigenvector \(u_n\in\mathscr C_{N_n,T_n}\) at a minimizing coordinate
mode.  With
\[
 \mathcal K_n=\mathscr C_{N_n,T_n},\qquad
 \mathfrak q_n[u]
 =g_{\mathrm{YM}}^{-2}\lVert\mathrm d_{N_n}u\rVert^2,
\]
\cref{cusp-gauge-finite-gap-limits} gives
\(\mathfrak q_n[u_n]=\gamma_{N_n,T_n}\to0\).  Thus the auxiliary
soft-state data in \cref{cusp-gauge-auxiliary-soft-states} are present.

What is absent is one already constructed, gauge-invariant,
positive-metric physical Hilbert space for four-dimensional quantum
Yang--Mills theory, its Hamiltonian \(H\), and maps \(J_n\) satisfying
\cref{cusp-gauge-bridge-norm,cusp-gauge-bridge-energy} with constants
controlled in physical units after gauge reduction, reconstruction,
renormalization, and all limiting operations.  Consequently
\cref{cusp-gauge-no-gap-bridge} cannot be applied to infer a quantum mass
statement from the cusp geometry.  Even if such maps were supplied, the
proposition would establish only spectrum arbitrarily close to zero for
that specified \(H\); it would not by itself identify a one-particle state,
establish scattering, or construct the required quantum gauge theory.
\end{remark}
